\documentclass{article}
\usepackage[T1]{fontenc}
\usepackage{lmodern}
\usepackage{graphicx}
\usepackage{amsmath,amssymb}
\usepackage[a4paper,margin=2cm]{geometry}
\usepackage[absolute,overlay]{textpos}
\setlength{\TPHorizModule}{1mm}
\setlength{\TPVertModule}{1mm}
\usepackage[indonesian]{babel}
\usepackage{hyperref}
\setlength{\emergencystretch}{2em}
% unit: Halaman 63

\begin{document}

% === Halaman 63 (python/native) ===

63

• GSM merupakan standard telepon seluler Eropa. A5 dikembangkan 

oleh Perancis

• Tidak semua operator GSM mengimplementasikan  enkripsi, 

bergantung regulasi (seperti di Indonesia)

• A5 ada dua versi: 


1. A5/1 : versi kuat A5, digunakan di Eropa


2. A5/2 (Kasumi) : versi ekspor, lebih lemah

• Algoritma A5/1 pada awalnya rahasia, tetapi pada tahun 1994 

melalui reverse engineering, algoritmanya terbongkar.

\end{document}
